\documentclass{beamer}
\usetheme{Madrid}

\usepackage[utf8]{inputenc}
\usepackage[T1]{fontenc}
\usepackage{amsmath, amssymb}
\usepackage{graphicx}
\usepackage{booktabs}

% To display speaker notes on a second screen, uncomment both lines below:
% \usepackage{pgfpages}
% \setbeameroption{show notes on second screen=right}

\title[CNN vs Transformer]{Comparing CNN and Transformer Architectures\\for Image Classification}
\subtitle{Advanced Deep Learning --- Project Presentation}
\author[B.~D.~Kir\'{a}ly]{B\'{a}lint D\'{a}niel Kir\'{a}ly}
\institute{E\"{o}tv\"{o}s Lor\'{a}nd University}
\date{June 2026}

% Fix footer: author | short title | date | frame/total
\setbeamertemplate{footline}{%
  \leavevmode%
  \hbox{%
    \begin{beamercolorbox}[wd=.333\paperwidth,ht=2.25ex,dp=1ex,center]{author in head/foot}%
      \usebeamerfont{author in head/foot}\insertshortauthor
    \end{beamercolorbox}%
    \begin{beamercolorbox}[wd=.334\paperwidth,ht=2.25ex,dp=1ex,center]{title in head/foot}%
      \usebeamerfont{title in head/foot}CNN vs Transformer
    \end{beamercolorbox}%
    \begin{beamercolorbox}[wd=.166\paperwidth,ht=2.25ex,dp=1ex,center]{date in head/foot}%
      \usebeamerfont{date in head/foot}\insertshortdate
    \end{beamercolorbox}%
    \begin{beamercolorbox}[wd=.166\paperwidth,ht=2.25ex,dp=1ex,right]{date in head/foot}%
      \usebeamerfont{date in head/foot}%
      \insertframenumber{} / \inserttotalframenumber\hspace*{2ex}
    \end{beamercolorbox}%
  }%
  \vskip0pt%
}

\begin{document}

% ---------------------------------------------------------------
\begin{frame}
  \titlepage
\end{frame}

% ---------------------------------------------------------------
\begin{frame}{The Task}

  \textbf{Objective:} Compare multiple CNN and Transformer models on image classification.

  \bigskip
  \textbf{Models:}
  \begin{itemize}
    \item Small CNN (trained from scratch)
    \item ResNet-18 (pretrained on ImageNet)
    \item ViT-Tiny (pretrained on ImageNet)
  \end{itemize}

  \bigskip
  \textbf{Datasets:} CIFAR-10 (10 classes) and CIFAR-100 (100 classes)

  \bigskip
  \textbf{Variables:} training strategy, classifier-head depth, and hyperparameters;
  comparison on multiple metrics.

  \note{This was an open-ended ``compare CNNs and Transformers'' brief.
  I turned it into a controlled set of ablations --- varying training strategy,
  head depth, and hyperparameters --- so that differences between models could be
  attributed to the right cause rather than confounded by the training recipe.}

\end{frame}

% ---------------------------------------------------------------
\begin{frame}{Experimental Setup}

  \begin{columns}
    \column{0.55\textwidth}

      \textbf{Models:}

      \medskip
      \small
      \renewcommand{\arraystretch}{1.3}
      \begin{tabular}{@{}lllr@{}}
        \toprule
        Model      & Type        & Pretrained & Params \\
        \midrule
        small\_cnn & CNN         & No  & 141{,}354 \\
        ResNet-18  & CNN         & Yes & 11{,}181{,}642 \\
        ViT-Tiny   & Transformer & Yes & 5{,}489{,}482 \\
        \bottomrule
      \end{tabular}

    \column{0.42\textwidth}

      \textbf{Protocol:}
      \begin{itemize}
        \item Early stopping (patience~5, max~20 epochs) with best-weight restoration
        \item 3 seeds (42\,/\,123\,/\,7) for the core comparison --- error bars confirm the ranking is real
        \item All cross-model comparisons locked to \emph{full fine-tuning + linear head}: differences reflect the \emph{model}, not the recipe
      \end{itemize}

  \end{columns}

  \note{Emphasise the fairness design. Early stopping gives every model the same
  convergence opportunity, and best-weight restoration ensures the reported number
  is the actual peak rather than a later overfit checkpoint. Three seeds show the
  ranking is stable, not a lucky draw. The fixed training condition is essential
  for attribution: you cannot claim one architecture beats another if the recipes differ.}

\end{frame}

% ---------------------------------------------------------------
\begin{frame}{Metrics}

  \textbf{Accuracy:} Top-1 and Top-5 validation accuracy

  \bigskip
  \textbf{Loss:} Best validation loss (cross-entropy)

  \bigskip
  \textbf{Model cost:} Trainable vs.\ total parameters;
  time per epoch; throughput (samples/sec)

  \bigskip
  \textbf{Generalisation gap:} Training Top-1 $-$ Validation Top-1

  \bigskip
  \textit{Note:} Top-5 is uninformative on CIFAR-10 (10 classes
  $\Rightarrow$ $\approx$99\,\% for all models) but becomes meaningful on CIFAR-100.

  \note{Explain why we report Top-5 but do not emphasise it on CIFAR-10: with only ten
  classes the correct answer is almost always in every model's top five, so the metric
  does not discriminate. On CIFAR-100 it becomes a useful secondary signal that separates
  models more clearly than Top-1 alone.}

\end{frame}

% ---------------------------------------------------------------
\begin{frame}{Core Comparison: CIFAR-10}

  \begin{center}
    \includegraphics[width=0.78\textwidth]{results/plots/core_comparison_cifar10.png}
  \end{center}

  \textbf{Under identical full fine-tuning} (mean\,$\pm$\,std over 3 seeds):
  \begin{itemize}
    \item ResNet-18: \textbf{85.5\,\%} \quad Small CNN: 81.3\,\% \quad ViT-Tiny: 69.4\,\%
    \item Tight error bars --- the ranking is real, not seed noise
  \end{itemize}

  \note{This is the headline result. Announce it clearly, but then immediately flag that
  ViT's low number reflects \emph{training difficulty}, not architectural capacity.
  The next two slides provide the evidence. Do not let this be read as ``CNNs beat
  Transformers'' until the audience has seen the full picture.}

\end{frame}

% ---------------------------------------------------------------
\begin{frame}{Training Dynamics}

  \begin{center}
    \includegraphics[width=0.80\textwidth]{results/plots/convergence_curves_cifar10.png}
  \end{center}

  \begin{itemize}
    \item ResNet-18 converges fastest and reaches the lowest validation loss
    \item ViT-Tiny is slowest and plateaus at the highest loss
    \item All runs are clearly converged --- the comparison is valid, not under-trained
  \end{itemize}

  \note{Use this as the ``results are trustworthy'' evidence. Every model ran to
  convergence under early stopping with best-weight restoration, so the accuracy figures
  are not underestimates from stopping too early. The curves also give the visual
  intuition that ViT is harder to optimise on this dataset, which sets up the key
  finding on the following slides.}

\end{frame}

% ---------------------------------------------------------------
\begin{frame}{Training Strategy Matters}

  \begin{center}
    \includegraphics[width=0.80\textwidth]{results/plots/strategy_comparison_cifar10.png}
  \end{center}

  \textbf{Pretrained backbones, CIFAR-10:}
  \begin{itemize}
    \item ResNet-18: full 85.5\,\% \;/\; freeze$\!\to\!$unfreeze 80.7\,\% \;/\; head-only 41.5\,\%
    \item ViT-Tiny:\phantom{xx} full 71.4\,\% \;/\; freeze$\!\to\!$unfreeze 79.7\,\% \;/\; head-only 40.0\,\%
    \item Head-only $\approx$40\,\% for both: frozen ImageNet features do not transfer to 32$\times$32 CIFAR images
    \item ViT \emph{gains} from staged unfreezing; ResNet does not need it
  \end{itemize}

  \note{The head-only failure is the same conceptual point for both models: the scale
  mismatch between ImageNet (224$\times$224) and CIFAR (32$\times$32) means frozen backbone
  features are wrong-scale. More interesting is that ViT benefits substantially from staged
  unfreezing while ResNet does not --- an early signal that ViT's optimisation is more
  sensitive to how training is structured.}

\end{frame}

% ---------------------------------------------------------------
\begin{frame}{Key Finding: ViT's Gap Is a Trainability Artifact}

  \begin{center}
    \includegraphics[width=0.78\textwidth]{results/plots/vit_lr_sensitivity_cifar10.png}
  \end{center}

  \textbf{ViT-Tiny full fine-tuning, varying learning rate:}
  \begin{itemize}
    \item lr $10^{-3}$ $\to$ 71.4\,\% \quad lr $3{\times}10^{-4}$ $\to$ 77.8\,\% \quad lr $10^{-4}$ $\to$ 79.7\,\%
    \item A gentler lr closes $\approx$8 points, nearly matching the CNNs
    \item Freeze$\!\to\!$unfreeze independently recovers to the same $\approx$80\,\%
  \end{itemize}

  \textit{Conclusion: ViT-Tiny is highly sensitive to fine-tuning choices;
  lr $10^{-3}$ is too aggressive. Do not read the core result as
  ``CNNs beat Transformers.''}

  \note{This is the intellectually honest centrepiece of the talk. Two entirely independent
  lines of evidence --- reducing the learning rate, and using staged unfreezing ---
  both converge on ViT reaching $\approx$80\,\%, closing most of the gap to ResNet.
  The weakness is in optimisation strategy, not in architectural capacity. Make this
  point explicitly before moving on.}

\end{frame}

% ---------------------------------------------------------------
\begin{frame}{Classifier Head Depth}

  \begin{center}
    \includegraphics[width=0.80\textwidth]{results/plots/head_depth_cifar10.png}
  \end{center}

  \begin{itemize}
    \item Linear / one-hidden / two-hidden differ by $\lesssim$1 point for small CNN and ResNet-18
    \item Deeper heads \emph{hurt} ViT-Tiny: 71.4\,\% $\to$ 66.8\,\%
    \item A linear head is sufficient; extra head capacity adds overfitting risk
  \end{itemize}

  \note{Quick slide with a clean negative result. Adding hidden layers to the classifier
  head does not help, and for ViT it actively hurts --- likely because the deeper head
  amplifies the same optimisation instability we saw in the learning-rate slide.
  The practical default should be a single linear classifier on top of the backbone.}

\end{frame}

% ---------------------------------------------------------------
\begin{frame}{Hyperparameter Sensitivity}

  \begin{center}
    \includegraphics[width=0.80\textwidth]{results/plots/hp_sensitivity_resnet18.png}
  \end{center}

  \textbf{ResNet-18 freeze$\!\to\!$unfreeze, varying lr and batch size:}
  \begin{itemize}
    \item Learning rate dominates: 60\,\% $\to$ 72\,\% $\to$ 81\,\%
          across lr $10^{-4}$ / $3{\times}10^{-4}$ / $10^{-3}$
    \item Batch size 64 vs.\ 128 is a minor effect at each lr level
    \item \textbf{Takeaway:} spend the tuning budget on learning rate
  \end{itemize}

  \note{This connects directly to the ViT learning-rate slide: across all models, lr is
  the dominant knob. Batch size matters far less. If you have a limited tuning budget for
  any image-classification transfer-learning problem, the learning rate is always the first
  and most important thing to tune.}

\end{frame}

% ---------------------------------------------------------------
\begin{frame}{Scaling to a Harder Dataset: CIFAR-100}

  \begin{columns}
    \column{0.50\textwidth}
      \includegraphics[width=\textwidth]{results/plots/dataset_scaling.png}
    \column{0.48\textwidth}
      \includegraphics[width=\textwidth]{results/plots/top5_cifar100.png}
  \end{columns}

  \bigskip
  \textbf{CIFAR-10 $\to$ CIFAR-100:}
  ResNet 85.5$\to$58.6 \,/\, small CNN 82.2$\to$55.0 \,/\, ViT 71.4$\to$43.2

  \begin{itemize}
    \item All drop $\approx$25--30 pts; \textbf{ordering is preserved}
    \item Top-5 on CIFAR-100 (84.5\,/\,83.2\,/\,72.8) is now informative
    \item Large models overfit on CIFAR-100 (gen.\ gap: ResNet $+18.6$, ViT $+16.7$) --- stronger regularisation is future work
    \item small\_cnn was still improving at 60 epochs; 55.0\,\% slightly underestimates its asymptote (ordering unaffected)
  \end{itemize}

  \note{Deliver the preserved ordering as the main robust claim --- it holds even on a
  much harder 100-class task. Then cover the honest limitations: large pretrained models
  overfit on CIFAR-100, suggesting regularisation or augmentation would help, and the
  small CNN had not quite converged. Neither caveat changes which model is best; they are
  honest notes for future work.}

\end{frame}

% ---------------------------------------------------------------
\begin{frame}{Efficiency: Accuracy vs.\ Cost}

  \begin{columns}
    \column{0.50\textwidth}
      \includegraphics[width=\textwidth]{results/plots/params_vs_accuracy_cifar10.png}
    \column{0.48\textwidth}
      \includegraphics[width=\textwidth]{results/plots/average_time_per_epoch_by_model.png}
  \end{columns}

  \bigskip
  \begin{itemize}
    \item \textbf{small CNN:} $\sim$0.14\,M params, $\sim$25\,s/epoch, 81\,\% --- remarkably efficient
    \item \textbf{ResNet-18:} best accuracy (85.5\,\%) at 11.2\,M params / $\sim$49\,s/epoch
    \item \textbf{ViT-Tiny:} more params and time than small CNN for less accuracy
          under na\"ive training (see ViT trainability slide for the full picture)
  \end{itemize}

  \note{Frame this as an accuracy-versus-cost trade-off. The small CNN is the efficiency
  surprise of the project: for a dataset of this scale a tiny well-designed CNN punches
  far above its weight. ResNet wins on accuracy but at 80 times the parameters. ViT's
  cost-effectiveness improves substantially once the learning rate is tuned --- so judge
  ViT under its best setting, not the naive default.}

\end{frame}

% ---------------------------------------------------------------
\begin{frame}{Conclusions}

  \begin{itemize}
    \item Under identical full fine-tuning,
          \textbf{ResNet-18 $>$ small CNN $>$ ViT-Tiny} on CIFAR-10
    \item ViT's gap is largely a \textbf{trainability artifact}: a gentle lr or
          staged unfreezing closes most of it
    \item \textbf{Transfer strategy matters} and is model-dependent;
          frozen ImageNet features alone are poor on CIFAR
    \item Head depth barely matters;
          \textbf{learning rate is the dominant hyperparameter}
    \item The harder dataset (CIFAR-100) \textbf{preserves the ranking};
          large models overfit there
    \item \textbf{Methodologically:} early stopping + best-weight restore +
          multi-seed error bars make the comparison fair and trustworthy
  \end{itemize}

  \note{Spend time on this slide. Go through each bullet as one spoken sentence: state
  the finding, then its significance. Linger on the methodological point last --- without
  those controls (early stopping, best weights, multiple seeds) the comparison would not
  be convincing, and those were deliberate design decisions, not defaults.}

\end{frame}

% ---------------------------------------------------------------
\begin{frame}{Use of Artificial Intelligence}

  \textbf{Transparency Statement:}

  AI tools were used in the following aspects of this project:
  \begin{itemize}
    \item \textbf{Code and pipeline:} the training and evaluation pipeline was
          built iteratively with AI assistance; I reviewed every change, diagnosed
          failures, and directed each round of redesign
    \item \textbf{Presentation:} these slides were generated by AI based on my
          content and requirements
    \item \textbf{Documentation:} AI produced initial drafts; all were
          substantially revised before final use
  \end{itemize}

  \bigskip
  \textit{I designed the experiments, identified the under-training issue,
  directed the full rerun, and verified every result.
  All interpretation is my own.}

\end{frame}

% ---------------------------------------------------------------
\begin{frame}{Thank You!}

  \begin{center}
    \Large Questions?
  \end{center}

  \bigskip
  \textbf{References:}
  \begin{itemize}
    \item Krizhevsky, A.\ (2009). \textit{Learning Multiple Layers of Features from
          Tiny Images}. (CIFAR-10/100 datasets)
    \item He, K., Zhang, X., Ren, S., \& Sun, J.\ (2016).
          \textit{Deep Residual Learning for Image Recognition}. CVPR.
    \item Dosovitskiy, A.\ et al.\ (2021). \textit{An Image is Worth 16$\times$16
          Words: Transformers for Image Recognition at Scale}. ICLR.
    \item Wightman, R.\ \textit{PyTorch Image Models} (\texttt{timm}).
          \texttt{github.com/rwightman/pytorch-image-models}
  \end{itemize}

\end{frame}

\end{document}
